\subsection{Operators}
\label{sec:operational-semantics-operators}

The comparisons \pyinline{==}, \pyinline{!=}, \pyinline{in} and \pyinline{not in} are given in terms of
structural equality $\eqOp{\Sigma}$ and membership $\contains{\Sigma}$. When defined, the
operators agree with Python, yielding $\aborts{\kappa}$ when Python raises an
exception. When undefined\paperonly{ (\secref{no-coercions})}, evaluation of the expression has no derivation, and an
implementation may abort or produce the result that Python gives (\paperonly{\secref{python-compatibility}}\speconly{\secref{conformance}}). Because
equality is partial, the order in which elements are compared is observable: equality and membership short-circuit
at the first element that decides the result, so an undefined comparison later in a sequence does not make
the whole comparison undefined.


Structural equality $\eqOp{\Sigma}$ compares literals by value, lists and tuples elementwise, dictionaries with
the same keys by their values, and objects of the same class by their fields. Values of different forms, such
as a number and a string or a list and a tuple, are unequal. Equality is undefined between a Boolean and a
number, which Python compares numerically, between functions, which Python compares by identity, and between
two \pyinline{nan}s inside a container. Compared directly, \pyinline{nan == nan} is \pyinline{False} in both languages,
but Python compares container elements by identity before equality, so two lists holding \pyinline{nan}s
are equal or not according to whether they hold the same object, a distinction that \PurePy cannot make.
Python gives an answer in each case (\figref{coercions}), so each is a dynamically excluded comparison. Comparison short-circuits at the first element that
decides the result, so the comparison order is observable:

\begin{python}
x: float = 1e400 - 1e400  # nan
print([1.0, x] == [2.0, x])  # False, decided at the first element
print([1.0, x] == [1.0, x])  # undefined: nan against nan
\end{python}

\noindent Dictionaries are compared in the key order of the left operand for the same reason. The membership
function $\contains{\Sigma}$ tests for a key of a dictionary, a substring of a string, or an element of a list or tuple,
the last by structural equality with the same short-circuiting. The strict operators are given by
$\hat{\odot}$ and $\hat{\ominus}$ (\figref{operators}); the definitions of all of these are given in
\specification.

\begin{figure}
\small
\begin{center}
\begin{tabular}{>{\raggedright\arraybackslash}p{0.22\textwidth}>{\raggedright\arraybackslash}p{0.62\textwidth}}
\arrayrulecolor{gray!60}\hline
\rowcolor{grammarheadcolor}$\odot$ & $\hat{\odot}(v, v')$ \\
\pyinline{==} & $\val{\pyinline{False}}$ if both operands are \pyinline{nan}, and $\eqOp{\Sigma}(v, v')$ otherwise \\
\pyinline{!=} & $\val{\pyinline{True}}$ if both operands are \pyinline{nan} or $\eqOp{\Sigma}(v, v') = \val{\pyinline{False}}$, $\val{\pyinline{False}}$ if $\eqOp{\Sigma}(v, v') = \val{\pyinline{True}}$, and otherwise $\eqOp{\Sigma}(v, v')$ \\
\pyinline{in} & $\contains{\Sigma}(v', v)$ \\
\pyinline{not in} & $\val{\pyinline{True}}$ if $\contains{\Sigma}(v', v) = \val{\pyinline{False}}$, $\val{\pyinline{False}}$ if $\contains{\Sigma}(v', v) = \val{\pyinline{True}}$, and otherwise $\contains{\Sigma}(v', v)$ \\
\pyinline{<}\quad\pyinline{<=}\quad\pyinline{>}\quad\pyinline{>=} & \todo{Ordering, as in Python} \\
\pyinline{+}\quad\pyinline{-}\quad\pyinline{*}\quad\pyinline{/}\quad\pyinline{//}\quad\pyinline{\%}\quad\pyinline{**} & \todo{Arithmetic, as in Python. Float precision is left to the implementation. $\aborts{\errZeroDiv}$ where the divisor is zero, and $\aborts{\errType}$ where an operand is not a number}
\end{tabular}


\vspace{2mm}

\begin{tabular}{>{\raggedright\arraybackslash}p{0.22\textwidth}>{\raggedright\arraybackslash}p{0.62\textwidth}}
\arrayrulecolor{gray!60}\hline
\rowcolor{grammarheadcolor}$\ominus$ & $\hat{\ominus}(v)$ \\
\pyinline{not} & $\val{\pyinline{True}}$ if $v = \pyinline{False}$, and $\val{\pyinline{False}}$ if $v = \pyinline{True}$ \\
\pyinline{+} & $\val{v}$ if $v$ is a number, and $\aborts{\errType}$ otherwise \\
\pyinline{-} & $\val{-v}$ if $v$ is a number, and $\aborts{\errType}$ otherwise
\end{tabular}

\end{center}
\caption{Operator meanings: value, abort, or undefined if no case applies}
\label{fig:operators}
\end{figure}

